\documentclass[journal]{IEEEtran}

\usepackage{cite}
\usepackage{graphicx}
\usepackage{amsmath,amssymb,amsfonts}
\usepackage{bm}
\usepackage{multirow}
\usepackage{array}
\usepackage{booktabs}
\usepackage{url}
\usepackage{xcolor}

\newcommand{\placeholderfig}[2]{%
    \fbox{\parbox[c][#2][c]{#1}{\centering\vspace{0.5em}Figure placeholder\\[0.5em]Final artwork pending\vspace{0.5em}}}%
}
\newcommand{\draftfigure}[3]{%
    \IfFileExists{#1}{\includegraphics[width=#2]{#1}}{\placeholderfig{#2}{#3}}%
}

\begin{document}

\title{HBTXR: Algorithm-Hardware Co-Designed Hybrid Eye Tracking for On-Device Extended Reality}

\author{Anonymous~Author(s)}

\markboth{IEEE Journal on Emerging and Selected Topics in Circuits and Systems,~Vol.~XX, No.~XX, Month~2026}%
{Anonymous \MakeLowercase{\textit{et al.}}: HBTXR}

\maketitle

\begin{abstract}
Eye tracking is a key enabling technology for extended reality (XR), supporting gaze-based interaction, foveated rendering, and privacy-preserving on-device perception~\cite{plopski2022eye,kwon2023xrbench}. However, existing approaches still suffer from a structural trade-off between semantic robustness and temporal responsiveness. Frame-based methods provide stable visual cues but incur high bandwidth cost and relatively large latency~\cite{li2018etracker,chen2022real}, whereas event-based methods offer high temporal resolution but remain limited in relocalization stability and representation consistency~\cite{li2024gaze,li2023track,ding2025facet}. To address these limitations, this paper presents \textbf{HBTXR}, an algorithm-hardware co-designed hybrid eye tracking system that tightly integrates \emph{Frame Search} and \emph{Event Track} for on-device XR applications. HBTXR is built upon four key components: 1) a dense-annotation-aware canonical training contract that stabilizes geometric supervision across heterogeneous frame and event inputs, 2) a hybrid frame-event eye tracking transformer with decoupled Search, Track, and Mask heads, 3) a two-stage host-side Search/Track scheduler for runtime-adaptive execution, and 4) an FPGA accelerator architecture with a shared transformer engine, streaming-based feature reuse, and FSM-based runtime control. Experimental results show that HBTXR achieves a pupil-center error of \textbf{0.42 pixels}, a gaze error of \textbf{0.68$^\circ$}, and an end-to-end latency of \textbf{0.57 ms} in Track mode, while maintaining a Search success rate of \textbf{97.8\%}. These results demonstrate that HBTXR provides a practical solution for real-time on-device eye tracking in next-generation XR systems.
\end{abstract}

\begin{IEEEkeywords}
Eye tracking, extended reality, event camera, hybrid sensing, transformer accelerator, FPGA, on-device inference.
\end{IEEEkeywords}

\section{Introduction}

Eye tracking has become a key enabling component of extended reality (XR) systems, supporting gaze-based interaction, foveated rendering, attention-aware interfaces, and privacy-preserving on-device perception~\cite{plopski2022eye,kwon2023xrbench}. In practical head-mounted XR devices, however, eye tracking must satisfy several requirements simultaneously: it must be semantically robust enough to recover the pupil under appearance variation, temporally responsive enough to track rapid eye motion with very low latency, and efficient enough to operate within tight on-device power, bandwidth, and form-factor constraints. These requirements make practical XR eye tracking fundamentally different from offline or server-assisted gaze estimation.

Conventional frame-based eye-tracking methods provide stable spatial cues and reliable semantic information for pupil localization and gaze estimation~\cite{li2018etracker,chen2022real,jain2025flattrack}. Their main advantage lies in global robustness: frame observations preserve appearance structure that is helpful for relocalization, geometric fitting, and calibration-aware regression. However, dense image acquisition incurs substantial memory traffic and bandwidth overhead, and its effective temporal responsiveness is limited when rapid eye motion must be tracked continuously. In XR settings, these drawbacks translate into higher latency and reduced suitability for tightly integrated on-device deployment.

Event-based eye tracking has therefore emerged as a compelling alternative~\cite{wang2024event,chen2025event}. By asynchronously sensing brightness changes, event cameras naturally provide sparse high-temporal-resolution observations and are well suited for low-latency update~\cite{li2024gaze,li2023track,ding2025facet,zhang2024swift,chen20233et,feng2022real,lai2026efficient,lenero20113}. Prior studies have shown that event-based methods are especially effective under fast eye motion and low-redundancy runtime requirements~\cite{li2023track,ding2025facet,chen20233et,feng2022real,lai2026efficient}. Nevertheless, purely event-based pipelines remain weaker in semantic stabilization and global recovery. When event density becomes low or drift accumulates over time, event-only tracking may lose the anchor needed for reliable relocalization. As a result, frame-only and event-only methods expose a structural trade-off: the former is semantically stable but temporally expensive, whereas the latter is temporally efficient but less reliable for long-horizon recovery.

Hybrid event-frame eye tracking has also been explored to bridge this gap by combining frame-guided relocalization with event-guided update~\cite{chen2025ex}. This direction is promising, but most prior hybrid systems still treat multimodal fusion primarily as a model-level problem. In many cases, supervision design, runtime switching, and deployment structure remain only loosely coupled. For practical XR eye tracking, however, the distinction between relocalization and local update is not merely architectural; it is also operational. Search and Track exhibit different input contexts, state dependencies, robustness requirements, and latency characteristics. This suggests that hybrid eye tracking should be treated as a system-level search-and-update problem rather than as a single undifferentiated multimodal regression task.

This work addresses that gap through \textbf{HBTXR}, an algorithm-hardware co-designed hybrid eye-tracking system for on-device XR applications. The core idea of HBTXR is to explicitly separate \emph{Frame Search} and \emph{Event Track} within a unified hybrid tracking framework. Search is responsible for robust relocalization and anchor refresh using semantically stable frame cues, whereas Track performs low-latency residual update conditioned on the previous state and recent event observations. To support this separation, HBTXR introduces geometry-stable hybrid supervision, a shared-backbone hybrid transformer with decoupled prediction heads, a runtime-aware Search/Track scheduler, and a deployment architecture based on mode-aware mixed-granularity execution.

Fig.~\ref{fig:overview} shows the overall concept of HBTXR. The system combines canonicalized frame-event inputs, a shared transformer-oriented representation path, decoupled Search/Event/Track/Mask heads, and a host-controlled scheduler that coordinates refresh-oriented and reuse-oriented execution. This overall organization is motivated by the observation that practical XR eye tracking requires supervision, model structure, runtime policy, and deployment mapping to be designed jointly rather than separately.

The main contributions of this work are summarized as follows.
\begin{itemize}
\item We propose a \textbf{geometry-stable hybrid supervision framework} that aligns frame-based relocalization and event-based residual tracking under a unified canonical target representation.
\item We design a \textbf{hybrid Search-Track transformer} with a shared backbone and decoupled Search, Event, Track, and Mask heads for robust relocalization and low-latency update.
\item We develop a \textbf{runtime-aware Search/Track scheduler} that switches between refresh-oriented and state-persistent operating modes using confidence, similarity, event density, and eye-state signals.
\item We present a \textbf{mode-aware mixed-granularity deployment architecture} that maps HBTXR onto a shared accelerator substrate with feature reuse and low-overhead control for on-device XR execution.
\end{itemize}

Accordingly, HBTXR is designed as a hybrid Search-and-Track system whose supervision, model structure, runtime policy, and deployment path are considered jointly rather than separately. The remainder of this paper is organized as follows. Section~II reviews prior work on frame-based, event-based, and hybrid eye tracking, as well as deployment-oriented accelerator designs. Section~III presents the proposed HBTXR algorithm. Section~IV describes the proposed HBTXR accelerator. Section~V reports the experimental setup, comparisons, and ablation studies. Finally, Section~VI concludes the paper.

\begin{figure*}[!t]
\centering
\draftfigure{fig1_overview.pdf}{0.95\textwidth}{2.2in}
\caption{Overview of HBTXR. The proposed system combines geometry-stable hybrid supervision, a Search-Track transformer, a runtime-aware host-side scheduler, and a mode-aware mixed-granularity deployment architecture for on-device XR eye tracking.}
\label{fig:overview}
\end{figure*}

\section{Related Work}

\subsection{Frame-Based Eye Tracking}

Frame-based eye tracking has long been the dominant approach for gaze estimation in head-mounted and near-eye systems~\cite{li2018etracker,chen2022real,jain2025flattrack}. These methods typically rely on RGB or infrared images to infer the pupil center, iris geometry, or gaze direction from semantically rich appearance cues. Their main advantage lies in stable spatial information, which makes relocalization, geometric fitting, and calibration-aware regression relatively straightforward. In practical systems, this often translates into reliable global recovery under moderate pose and appearance variation.

However, frame-based pipelines remain limited in latency-critical XR scenarios. Because they depend on dense image acquisition and repeated full-frame processing, they incur substantial bandwidth, memory, and compute overhead. Their temporal responsiveness is also constrained when rapid eye motion must be tracked continuously. As a result, although frame-based methods provide strong semantic stability, they are not ideal as standalone solutions for low-latency on-device XR eye tracking.

\subsection{Event-Based Eye Tracking}

Event-based eye tracking exploits the asynchronous sensing property of event cameras to achieve low-latency and low-redundancy processing~\cite{li2024gaze,li2023track,ding2025facet,zhang2024swift,sen2024eyetraes,chen20233et,chen2025event,feng2022real,lai2026efficient,wang2024event,lenero20113,zhao2025ev}. Representative methods such as E-Track directly estimate pupil states from event representations~\cite{li2023track}, while FACET demonstrates that geometry-aware event modeling can yield efficient and accurate pupil localization under dense supervision~\cite{ding2025facet}. Other works improve robustness or efficiency through anti-blink design, sparse processing, or spiking computation~\cite{zhang2024swift,sen2024eyetraes,sui2026event}.

Despite these advantages, event-only pipelines remain weaker in semantic stabilization and re-initialization. When event support becomes weak or drift accumulates over time, event-only trackers may lose the global anchor needed for reliable relocalization. This makes event-based tracking appealing for short-horizon update, but less sufficient as a complete standalone solution for robust XR deployment.

\subsection{Hybrid Event-Frame Eye Tracking}

Hybrid eye tracking combines the semantic stability of frames with the temporal responsiveness of events~\cite{chen2025ex,zhao2023ev}. EX-Gaze is a representative example showing that frame-based relocalization and event-based tracking can be integrated to improve both robustness and responsiveness~\cite{chen2025ex}. The EV-Eye study further highlights the importance of jointly exploiting semantic frame cues and high-frequency event observations for near-eye tracking~\cite{zhao2023ev}.

While promising, most existing hybrid methods still focus primarily on model-level fusion. In many cases, supervision design, runtime mode switching, and deployment structure remain weakly coupled. For practical XR eye tracking, however, Search and Track exhibit different functional and operational roles and should therefore be treated explicitly rather than being absorbed into a single multimodal regressor.

\subsection{Deployment-Oriented and Hardware-Aware Eye Tracking}

A separate but increasingly important direction focuses on making eye tracking practical for embedded and XR devices through deployment-aware design and hardware acceleration~\cite{you2022eyecod,han2026janeeye,giordano2025sub,bonazzi2024retina}. EyeCoD demonstrates optics-algorithm co-design for compact eye-tracking systems~\cite{you2022eyecod}, whereas recent event-driven deployments such as Retina, Sub-Millisecond MCU, and JaneEye show that low-latency eye tracking can be realized on neuromorphic processors, microcontrollers, and ASIC-style hardware~\cite{bonazzi2024retina,giordano2025sub,han2026janeeye}. These studies confirm that low-latency eye tracking must be evaluated not only at the model level, but also at the execution and data-movement levels.

Recent accelerator-oriented studies in this area further suggest that execution efficiency depends strongly on memory access patterns, synchronization overhead, and workload orchestration in addition to arithmetic throughput. These observations are directly relevant to hybrid eye tracking, where refresh-oriented Search and reuse-oriented Track exhibit different execution characteristics. However, prior deployment-oriented eye-tracking systems have largely focused either on frame-based pipelines or on event-only sparse paths, leaving a gap for hybrid Search/Track systems with transformer-oriented execution structure.

\subsection{Position of HBTXR}

HBTXR is positioned at the intersection of these directions. Unlike frame-only methods, it does not rely on repeated dense-image inference for steady-state tracking. Unlike event-only methods, it does not assume that local event evidence is sufficient for stable long-horizon recovery. Unlike many prior hybrid approaches, it explicitly separates Search and Track as distinct operating modes. Finally, unlike prior deployment-oriented systems specialized to a single sensing path or a fixed inference schedule, HBTXR is designed as an algorithm-hardware co-designed hybrid tracking system in which supervision, branch decomposition, runtime scheduling, and deployment mapping are considered jointly.

This gap motivates HBTXR as a runtime-aware and deployment-aware hybrid eye-tracking system, whose algorithmic structure is described in Section~III and whose accelerator-oriented mapping is presented in Section~IV.

\section{The Proposed HBTXR Algorithm}

\subsection{Problem Formulation}

As illustrated in Fig.~\ref{fig:hbtxr_algorithm}, HBTXR is designed for on-device XR eye tracking under a hybrid sensing setting composed of a frame stream, an event stream, and a recurrent tracking state. Let \(I_t \in \mathbb{R}^{H \times W}\) denote the frame acquired at time \(t\), and let
\begin{equation}
\mathcal{E}_t = \left\{ e_k = (x_k, y_k, p_k, \tau_k) \right\}_{k=1}^{N_t}
\label{eq:event_stream}
\end{equation}
denote the event set collected over the current temporal window, where \((x_k,y_k)\) is the event location, \(p_k \in \{-1,+1\}\) is the polarity, and \(\tau_k\) is the timestamp. The event stream is converted into a polarity-split event tensor \(V_t \in \mathbb{R}^{2 \times H \times W}\) through a causal accumulation operator
\begin{equation}
V_t = \Phi(\mathcal{E}_t; \omega_t),
\label{eq:event_tensor}
\end{equation}
where \(\omega_t\) denotes the event-window construction policy such as fixed-count or time-bin slicing.

The objective of HBTXR is to estimate the current pupil state
\begin{equation}
\mathbf{s}_t = [x_t, y_t, a_t, b_t, u_t, v_t] \in \mathbb{R}^{6},
\label{eq:state_def}
\end{equation}
where \((x_t,y_t)\) is the pupil center, \((a_t,b_t)\) are the ellipse axes, and \((u_t,v_t)\) is a trigonometric orientation encoding. In particular, HBTXR adopts the relation
\begin{equation}
[u_t, v_t] = [\sin(2\theta_t), \cos(2\theta_t)],
\label{eq:uv_def}
\end{equation}
which avoids discontinuity around the angular wrap-around boundary and enables stable regression of ellipse orientation.

A central design choice of HBTXR is to formulate hybrid eye tracking as a Search-and-Track problem rather than as a single monolithic prediction task. The Search path is responsible for robust relocalization and anchor refresh using frame-based semantic cues, while the Track path performs low-latency residual update from the previous state using event observations. In addition, an event-conditioned estimation branch is introduced to provide an auxiliary event-side state hypothesis that supports cross-branch consistency and runtime confidence assessment.

\subsection{Geometry-Stable Hybrid Supervision}

A major challenge in hybrid frame-event tracking is that frame cues, event cues, and recurrent state updates must be trained against a common geometric target. To address this issue, HBTXR adopts a geometry-stable hybrid supervision framework. Each raw sample is converted into a canonical training representation containing an eye region, pupil mask, ellipse annotation, and quality metadata. Let \(\mathcal{T}(\cdot)\) denote the canonicalization operator. The training sample is constructed as
\begin{equation}
(\bar{I}_t, \bar{V}_t, \bar{\mathbf{s}}_{t-1}, \bar{\mathbf{s}}_t, \bar{M}_t, \bar{B}_t, \bar{\xi}_t)
=
\mathcal{T}(I_t, \mathcal{E}_t, \mathbf{s}_{t-1}, \mathbf{s}_t, M_t, B_t, \xi_t),
\label{eq:canonical_contract}
\end{equation}
where \(\bar{I}_t\) and \(\bar{V}_t\) denote the canonicalized frame and event inputs, \(\bar{M}_t\) is the canonical pupil mask, \(\bar{B}_t\) is the canonical eye-region target, and \(\bar{\xi}_t\) collects supervision-quality signals such as annotation confidence, mask validity, closed-eye flag, and track validity.

This supervision contract provides three important properties. First, Search and Track are trained against the same canonical state representation, which prevents the two branches from drifting toward incompatible geometries. Second, the eye region, mask, ellipse, and state targets are spatially aligned under a single transform, so the supervision remains consistent across modalities. Third, quality and validity flags are propagated explicitly to both the loss functions and the scheduler, allowing the model to distinguish between reliable tracking conditions and samples that should be partially masked during training.

\subsection{Hybrid Search-Track Transformer}

The proposed HBTXR network is shown in Fig.~\ref{fig:hbtxr_algorithm}. The model uses modality-specific embedding layers and adapters followed by a shared transformer backbone. Rather than fully fusing frame and event tensors at the input level, HBTXR preserves modality-specific front-ends and reuses a common token-processing trunk so that search-oriented frame inference and track-oriented event inference can be expressed under a unified backbone.

For the frame branch, the canonical frame \(\bar{I}_t\) is converted into patch tokens through a frame patch embedding module:
\begin{equation}
\mathbf{Z}^{f}_t = \mathcal{B}\!\left(\mathcal{A}_f\!\left(\mathcal{P}_f(\bar{I}_t)\right)\right),
\label{eq:frame_embed}
\end{equation}
where \(\mathcal{P}_f(\cdot)\) denotes frame patch embedding, \(\mathcal{A}_f(\cdot)\) is a modality adapter, and \(\mathcal{B}(\cdot)\) is the shared Partial DeiT-Tiny backbone. For the event branch, the canonical event tensor \(\bar{V}_t\) is processed similarly:
\begin{equation}
\mathbf{Z}^{e}_t = \mathcal{B}\!\left(\mathcal{A}_e\!\left(\mathcal{P}_e(\bar{V}_t)\right)\right).
\label{eq:event_embed}
\end{equation}

On top of the shared backbone, HBTXR uses decoupled prediction heads: an Eye Region Head, a Pupil Search Head, an Event Search Head, a Pupil Track Head, a Search Mask Head, and an Auxiliary State Head. Let \(\mathbf{p}^{f}_t\) and \(\mathbf{p}^{e}_t\) denote the pooled frame and event features, respectively. The search branch predicts
\begin{equation}
\hat{\mathbf{y}}^{s}_t = h_s(\mathbf{p}^{f}_t),
\qquad
\hat{\mathbf{s}}^{s}_t = \hat{\mathbf{y}}^{s}_t[1{:}6],
\label{eq:search_pred}
\end{equation}
while the event branch predicts
\begin{equation}
\hat{\mathbf{y}}^{e}_t = h_e(\mathbf{p}^{e}_t),
\qquad
\hat{\mathbf{s}}^{e}_t = \hat{\mathbf{y}}^{e}_t[1{:}6].
\label{eq:event_pred}
\end{equation}
The dense mask branch predicts
\begin{equation}
\hat{M}_t = h_m(\mathbf{Z}^{f}_t),
\label{eq:mask_pred}
\end{equation}
which provides both geometric regularization and a confidence-aware auxiliary signal for relocalization quality.

\begin{figure*}[!t]
\centering
\draftfigure{fig2_algorithm.pdf}{0.95\textwidth}{2.2in}
\caption{Architecture of the proposed HBTXR algorithm. Frame and event inputs are processed by modality-specific embedding paths and a shared transformer backbone, followed by decoupled Search, Event, Track, and Mask heads.}
\label{fig:hbtxr_algorithm}
\end{figure*}

\subsection{Search Branch, Track Branch, and State Update}

The Search branch is designed for global relocalization. It is frame-guided because frame observations are more semantically stable under challenging event conditions. In contrast, the Track branch is designed for high-rate residual update. It combines the event feature \(\mathbf{p}^{e}_t\) with an encoded previous state \(\psi(\bar{\mathbf{s}}_{t-1})\) and predicts a residual vector:
\begin{equation}
\Delta \hat{\mathbf{r}}_t
=
h_t\!\left(
\left[
\mathbf{p}^{e}_t ; \psi(\bar{\mathbf{s}}_{t-1})
\right]
\right),
\label{eq:track_residual}
\end{equation}
where
\begin{equation}
\Delta \hat{\mathbf{r}}_t =
[\Delta \hat{x}_t, \Delta \hat{y}_t, \Delta \widehat{\log a}_t, \Delta \widehat{\log b}_t, \Delta \hat{u}_t, \Delta \hat{v}_t, \hat{c}^{trk}_t, \hat{q}^{trk}_t].
\label{eq:track_output}
\end{equation}
Here, \(\hat{c}^{trk}_t\) denotes track confidence and \(\hat{q}^{trk}_t\) denotes track quality.

The decoded tracking state is obtained as
\begin{equation}
\hat{x}_t = x_{t-1} + \Delta \hat{x}_t, \qquad
\hat{y}_t = y_{t-1} + \Delta \hat{y}_t,
\label{eq:track_xy}
\end{equation}
\begin{equation}
\hat{a}_t = a_{t-1}\exp(\Delta \widehat{\log a}_t), \qquad
\hat{b}_t = b_{t-1}\exp(\Delta \widehat{\log b}_t),
\label{eq:track_ab}
\end{equation}
\begin{equation}
[\hat{u}_t,\hat{v}_t]
=
\frac{
[u_{t-1} + \Delta \hat{u}_t,\; v_{t-1} + \Delta \hat{v}_t]
}{
\left\|
[u_{t-1} + \Delta \hat{u}_t,\; v_{t-1} + \Delta \hat{v}_t]
\right\|_2 + \epsilon
}.
\label{eq:track_uv}
\end{equation}
This residual parameterization is particularly suitable for short-horizon event-driven updates because it models tracking as a constrained local update around the current anchor.

\subsection{Runtime-Aware Search/Track Scheduler}

A key component of HBTXR is the explicit Search/Track scheduler. Rather than using all branches uniformly at every step, the scheduler decides whether the system should remain in low-latency Track mode or fall back to Search mode for relocalization. This decision is based on search confidence, track confidence, track quality, cross-branch similarity, event density, and eye-state validity. Let \(c^s_t\) be the search confidence, \(c^{trk}_t\) the tracking confidence, \(q^{trk}_t\) the tracking quality, \(\rho_t\) the similarity between search and track states, \(d_t\) the event density, and \(o_t\) a binary closed-eye indicator. The scheduler transitions from Search to Track when the search estimate is reliable:
\begin{equation}
\text{Search} \rightarrow \text{Track}
\quad \text{if} \quad
c^s_t > \tau_s.
\label{eq:scheduler_st}
\end{equation}
Conversely, it re-enters Search mode if the tracking path becomes unreliable:
\begin{equation}
\text{Track} \rightarrow \text{Search}
\quad \text{if} \quad
\left(c^{trk}_t \le \tau_t\right)
\vee
\left(q^{trk}_t \le \tau_q\right)
\vee
\left(\rho_t \le \tau_{\rho}\right)
\vee
\left(d_t \le \tau_d\right)
\vee
(o_t = 1).
\label{eq:scheduler_ts}
\end{equation}

This scheduler serves two roles. Algorithmically, it enforces the distinction between global relocalization and local residual update. Operationally, it exposes a runtime structure that can be mapped naturally to a host-controlled accelerator path. The scheduler is therefore not only a tracking policy, but also a key part of the algorithm-hardware co-design philosophy of HBTXR.

\begin{figure}[!t]
\centering
\draftfigure{fig3_scheduler.pdf}{0.95\columnwidth}{1.55in}
\caption{Runtime-aware Search/Track scheduler. The scheduler switches between refresh-oriented Search and state-persistent Track according to confidence, similarity, event density, and eye-state signals.}
\label{fig:scheduler}
\end{figure}

\subsection{Training Strategy and Loss Design}

HBTXR is trained in two stages. In the first stage, the model is trained primarily for Search-oriented relocalization so that the frame branch learns robust global geometric estimation. In the second stage, the hybrid network is fine-tuned with Track supervision, event-conditioned estimation, and cross-branch consistency. This stage-wise strategy is important because relocalization and residual tracking have different optimization characteristics and failure modes.

The total training objective is written as
\begin{equation}
\mathcal{L}
=
\lambda_{eye}\mathcal{L}_{eye}
+
\lambda_{mask}\mathcal{L}_{mask}
+
\lambda_{search}\mathcal{L}_{search}
+
\lambda_{event}\mathcal{L}_{event}
+
\lambda_{track}\mathcal{L}_{track}
+
\lambda_{cons}\mathcal{L}_{cons}
+
\lambda_{ctr}\mathcal{L}_{ctr}
+
\lambda_{aux}\mathcal{L}_{aux},
\label{eq:total_loss}
\end{equation}
where \(\mathcal{L}_{eye}\) is the eye-region loss, \(\mathcal{L}_{mask}\) is the pupil-mask loss, \(\mathcal{L}_{search}\) supervises frame-based search, \(\mathcal{L}_{event}\) supervises event-conditioned estimation, \(\mathcal{L}_{track}\) supervises residual tracking, \(\mathcal{L}_{cons}\) enforces cross-branch consistency, \(\mathcal{L}_{ctr}\) constrains the predicted center within the valid search region, and \(\mathcal{L}_{aux}\) denotes auxiliary supervision.

For both search and tracking, HBTXR uses geometry-aware supervision in addition to standard center and axis regression. Let \(\mu(\mathbf{s})\) and \(\Sigma(\mathbf{s})\) denote the ellipse center and covariance implied by state \(\mathbf{s}\). A geometry-aware term is defined as
\begin{equation}
\mathcal{L}_{geo}(\hat{\mathbf{s}}, \mathbf{s})
=
\left\|
\mu(\hat{\mathbf{s}})-\mu(\mathbf{s})
\right\|_2^2
+
\left\|
\Sigma(\hat{\mathbf{s}})^{1/2}-\Sigma(\mathbf{s})^{1/2}
\right\|_F^2,
\label{eq:geo_loss}
\end{equation}
which stabilizes the full ellipse geometry rather than supervising only the center location. Quality-aware loss gating is further applied based on annotation confidence, mask validity, tracking validity, and closed-eye status.

The algorithmic decomposition above directly informs the deployment structure of HBTXR, where refresh-oriented Search and state-persistent Track are mapped onto a shared accelerator substrate, as described in Section~IV.

\section{The Proposed HBTXR Accelerator}

\subsection{Design Challenges and Co-Design Principles}

The deployment objective of HBTXR is not merely to accelerate a generic transformer, but to support a hybrid eye-tracking workload whose operating modes exhibit fundamentally different execution characteristics. Search performs robust relocalization and anchor refresh from frame-guided features, whereas Track performs event-driven residual update conditioned on the previous state. This asymmetry leads to three hardware design challenges.

First, Search and Track have different dataflow requirements. Search is refresh-oriented and requires broader feature regeneration, whereas Track is reuse-oriented and benefits from persistent state and incremental update. A uniform inference path therefore leads either to unnecessary recomputation during Track mode or to insufficient global context during Search mode. Second, the shared transformer backbone introduces synchronization-heavy stages, especially in token interaction and branch fusion, while the decoupled heads and recurrent state update remain more amenable to lightweight streaming execution. Third, the target deployment setting imposes tight latency and memory constraints, making it undesirable to duplicate complete frame and event pipelines or to rely on repeated full-feature buffering.

These observations motivate an algorithm-hardware co-design strategy. In particular, the shared-backbone and decoupled-head design enables the hardware to reuse a common transformer-oriented compute substrate, while the explicit Search/Track scheduler exposes two runtime modes that can be mapped differently in terms of buffering, feature reuse, and execution granularity.

\subsection{CPU-FPGA Workload Partition}

Fig.~\ref{fig:accel_overview} illustrates the overall deployment architecture of HBTXR. The system adopts a heterogeneous CPU-FPGA organization in which the host processor is responsible for control-centric functions and the FPGA is responsible for latency-dominant neural inference. This partition is chosen because Search/Track scheduling, state management, and control-policy updates are irregular and mode-dependent, whereas token generation, transformer inference, and head execution are regular and amenable to spatial acceleration.

The CPU handles frame/event acquisition, runtime mode management, threshold comparison, closed-eye and validity handling, and final state dispatch. The FPGA handles frame/event tokenization, modality adaptation, shared-backbone inference, Search/Event/Track/Mask head execution, and mode-specific feature reuse. Search mode requires stronger global robustness and therefore benefits from explicit host-visible control and a refresh-oriented execution path. Track mode, in contrast, is dominated by short-horizon low-latency updates and is therefore mapped to a state-persistent inference path with minimal control overhead.

\subsection{Mode-Aware Mixed-Granularity Execution Architecture}

A key design choice of the HBTXR accelerator is to avoid a single uniform execution granularity for the entire hybrid tracking pipeline. Instead, the proposed accelerator adopts a \emph{mode-aware mixed-granularity execution architecture}, shown in Fig.~\ref{fig:accel_overview}. The main idea is that different stages of HBTXR exhibit different locality and synchronization characteristics and should therefore be mapped using different execution styles.

Stages that require broader synchronization across tokens or branches are executed using a \emph{phase-level execution path}. This path is appropriate for operations such as shared transformer processing, branch coordination, and full Search refresh, where reusable context is limited and inter-token dependency is relatively high. In contrast, stages that benefit from sequential locality and partial reuse are executed using a \emph{tile-stream execution path}. This path is suitable for Track-mode updates, lightweight head execution, and recurrent state propagation, where the dominant objective is to minimize idle gaps and buffering overhead while maintaining low-latency response.

\begin{figure*}[!t]
\centering
\draftfigure{fig4_accelerator.pdf}{0.95\textwidth}{2.2in}
\caption{Mode-aware mixed-granularity execution architecture of the proposed HBTXR accelerator. Phase-level execution is used for globally synchronized stages, whereas tile-stream execution is used for reuse-friendly update stages.}
\label{fig:accel_overview}
\end{figure*}

\subsection{Shared Backbone Engine and Mode-Specific Head Mapping}

The computational core of the accelerator is a shared transformer-oriented backbone engine that executes the common representation trunk of HBTXR. Rather than implementing separate hardware paths for frame Search and event Track, the accelerator maintains a common backbone execution substrate and attaches mode-specific heads on top of it.

Let \(T_f\) and \(T_e\) denote the frame and event token sequences. The shared representation engine computes
\begin{equation}
\mathbf{H}_f = \mathcal{B}(T_f), \qquad
\mathbf{H}_e = \mathcal{B}(T_e),
\label{eq:shared_backbone_accel}
\end{equation}
where \(\mathcal{B}(\cdot)\) is realized by a common transformer engine. In Search mode, the accelerator emphasizes frame-guided refresh and generates search-oriented outputs from the shared representation. In Track mode, the accelerator prioritizes event-conditioned update and combines event features with the encoded previous state. The Search head and Mask head are activated primarily during refresh-oriented operation, while the Track head dominates steady-state low-latency execution. This avoids duplicating the dominant compute blocks while preserving functional separation.

\subsection{Memory Organization and Reuse-Oriented Dataflow}

Because the target application is on-device XR eye tracking, memory movement is often more critical than arithmetic count. The proposed accelerator therefore distinguishes between \emph{refresh-oriented buffering} and \emph{stream-oriented reuse}. Search mode uses a refresh buffer pair to support broader feature regeneration and robust relocalization. Track mode uses a streaming token queue and persistent state buffers so that recently computed context can be reused without re-materializing the entire inference path.

This organization is shown conceptually in Fig.~\ref{fig:search_track_timeline}. During Search mode, frame and event inputs are staged into a refresh path that regenerates the required shared features. During Track mode, the accelerator reuses persistent features and updates only the event-conditioned path together with the previous-state encoding. In this way, the dominant cost of Search is paid only when anchor refresh is required, while Track remains lightweight and reuse-oriented.

\begin{figure}[!t]
\centering
\draftfigure{fig5_timeline.pdf}{0.95\columnwidth}{1.55in}
\caption{Search-path and Track-path execution timelines. Search follows a refresh-heavy path with stronger synchronization, whereas Track follows a reuse-oriented path with persistent state and reduced idle gaps.}
\label{fig:search_track_timeline}
\end{figure}

\subsection{Runtime Control and Scheduling Support}

To support mode-aware execution, the accelerator is coupled with an FSM-based runtime controller that interprets the host-side Search/Track decisions and dispatches the corresponding execution path. The controller manages three classes of signals: 1) mode-control signals from the host, 2) data-valid and buffer-status signals within the accelerator, and 3) branch-confidence and output-ready signals used to synchronize execution and state update. This design matches the scheduler semantics introduced in Section~III.

This control structure is one of the main distinctions between HBTXR and a conventional transformer accelerator. In a generic accelerator, the dominant concern is maximizing uniform utilization over a fixed computation graph. In HBTXR, however, the execution graph itself is mode-dependent. The controller therefore realizes the runtime semantics of Search and Track as two distinct operating states of the system.

\subsection{Latency Model and Deployment Scope}

The end-to-end latency of HBTXR can be decomposed as
\begin{equation}
T_{\mathrm{e2e}}
=
T_{\mathrm{host}}
+
T_{\mathrm{transfer}}
+
T_{\mathrm{exec}}
+
T_{\mathrm{sync}},
\label{eq:latency_decomp}
\end{equation}
where \(T_{\mathrm{host}}\) denotes host-side control and scheduling time, \(T_{\mathrm{transfer}}\) denotes host-device communication, \(T_{\mathrm{exec}}\) denotes accelerator execution time, and \(T_{\mathrm{sync}}\) denotes synchronization overhead. Because the Search and Track paths have different execution structures, their latencies are also different. Let \(T_{\mathrm{search}}\) and \(T_{\mathrm{track}}\) denote the two mode-specific latencies. The average runtime under scheduler-driven operation can be expressed as
\begin{equation}
T_{\mathrm{avg}}
=
\alpha T_{\mathrm{search}}
+
(1-\alpha)T_{\mathrm{track}},
\label{eq:avg_latency}
\end{equation}
where \(\alpha\) is the fraction of Search invocations.

The purpose of the proposed accelerator is therefore not to maximize a synthetic throughput benchmark in isolation, but to provide a deployment-oriented execution architecture that aligns with the actual Search/Track operating behavior of HBTXR. Based on the proposed algorithm and accelerator organization, Section~V evaluates HBTXR from the perspectives of tracking accuracy, runtime behavior, deployment-oriented efficiency, and component-wise ablation.

\section{Experimental Results}

\subsection{Experimental Setup}

\subsubsection{Dataset and Evaluation Protocol}

We evaluate HBTXR on a near-eye hybrid eye-tracking benchmark composed of synchronized frame observations, event streams, and dense pupil annotations. Each sample contains a canonicalized frame patch, a polarity-split event tensor, the previous tracking state, and geometry-aware supervision targets including the pupil mask, eye region, and current pupil state. All samples are transformed into a unified canonical coordinate system before training and evaluation.

For robust generalization assessment, the dataset is divided using a subject-disjoint protocol. The training split is used for stage-wise model optimization, the validation split is used for hyperparameter selection and early stopping, and the test split is used only for final reporting. The detailed configuration of the training protocol, event policy, runtime settings, and optimizer is summarized in Table~\ref{tab:exp_setup}.

\subsubsection{Training and Runtime Configuration}

HBTXR is trained in two stages. In the first stage, the model is optimized for Search-oriented relocalization so that the frame-guided branch learns stable eye-region and pupil-state estimation. In the second stage, the hybrid network is fine-tuned with Search, Event, and Track supervision jointly, together with consistency and quality-aware gating. As shown in Table~\ref{tab:exp_setup}, the default event representation uses a fixed-count window with 5000 events and fast causal accumulation, while all inputs are canonicalized to a direct-square \(256 \times 256\) representation.

At runtime, the Search/Track scheduler uses confidence, similarity, event-density, and eye-state signals to select between refresh-oriented Search and state-persistent Track. This evaluation setting is consistent with the deployment model described in Section~IV.

\subsubsection{Evaluation Metrics}

We evaluate the tracking accuracy using pupil-center error and gaze angular error. Let \(\hat{\mathbf{c}}_t\) and \(\mathbf{c}_t\) denote the predicted and ground-truth pupil centers at time \(t\). The average pupil-center error is defined as
\begin{equation}
\mathrm{Err}_{\mathrm{pupil}}
=
\frac{1}{N}
\sum_{t=1}^{N}
\left\|
\hat{\mathbf{c}}_t - \mathbf{c}_t
\right\|_2.
\label{eq:pupil_err}
\end{equation}
Similarly, if \(\hat{\mathbf{g}}_t\) and \(\mathbf{g}_t\) denote the predicted and ground-truth gaze vectors, the average gaze error is computed as
\begin{equation}
\mathrm{Err}_{\mathrm{gaze}}
=
\frac{1}{N}
\sum_{t=1}^{N}
\arccos
\left(
\frac{
\hat{\mathbf{g}}_t^\top \mathbf{g}_t
}{
\|\hat{\mathbf{g}}_t\|_2 \|\mathbf{g}_t\|_2
}
\right).
\label{eq:gaze_err}
\end{equation}
In addition, we report end-to-end latency, effective update rate, and Search success rate when applicable.

\subsubsection{Baselines}

We compare HBTXR with representative prior methods from three categories: 1) frame-based eye tracking, 2) event-based eye tracking, and 3) hybrid or deployment-oriented eye-tracking systems. The baseline categories are summarized in Table~\ref{tab:baseline_summary}.

\begin{table}[!t]
\caption{Experimental Setup Summary}
\label{tab:exp_setup}
\centering
\footnotesize
\begin{tabular}{ll}
\toprule
Item & Setting \\
\midrule
Training protocol & Two-stage Search pretrain + hybrid finetune \\
Input frame size & $256 \times 256$ \\
Input event size & $2 \times 256 \times 256$ \\
Event policy & Fixed-count \\
Default event count & 5000 \\
Accumulation & Fast causal accumulation \\
Resize policy & Direct square canonical resize \\
Scheduler signals & Search conf., track conf., similarity, density, eye state \\
Optimizer & AdamW \\
Batch size & 8 \\
Learning rate & $3\times10^{-4}$ \\
Weight decay & $1\times10^{-4}$ \\
Runtime platform & Host-device heterogeneous execution \\
\bottomrule
\end{tabular}
\end{table}

\begin{table}[!t]
\caption{Baseline Categories Used in Comparison}
\label{tab:baseline_summary}
\centering
\footnotesize
\begin{tabular}{lll}
\toprule
Category & Representative Methods & Primary Focus \\
\midrule
Frame-based & Etracker, HE-Tracker & Semantic stability \\
Event-based & E-Track, FACET, Swift-Eye & Low-latency update \\
Hybrid / deploy. & EX-Gaze, SEE, EyeCoD, JaneEye & Fusion / deployment \\
\bottomrule
\end{tabular}
\end{table}

\subsection{Main Results}

Table~\ref{tab:main_results} summarizes the main performance of HBTXR. In Search mode, HBTXR achieves a pupil-center error of 0.58 pixels, a gaze error of 0.91$^\circ$, a Search success rate of 97.8\%, and a latency of 0.93 ms. In Track mode, the proposed system achieves a pupil-center error of \textbf{0.42 pixels}, a gaze error of \textbf{0.68$^\circ$}, and an end-to-end latency of \textbf{0.57 ms}, corresponding to an effective update rate of \textbf{1754 Hz}. Under scheduler-driven operation, the average runtime is reduced to 0.63 ms while preserving a pupil-center error of 0.45 pixels and a gaze error of 0.71$^\circ$.

These results indicate that HBTXR maintains both geometric accuracy and sub-millisecond responsiveness during steady-state tracking, which is essential for practical XR applications. Moreover, the contrast between Search and Track confirms the central design intuition of HBTXR: Search should be used as a refresh-oriented recovery path, whereas Track should dominate low-latency steady-state operation.

Fig.~\ref{fig:qualitative_results} further visualizes representative Search and Track outputs. As highlighted in the caption of Fig.~\ref{fig:qualitative_results}, HBTXR maintains stable local updates in Track mode and re-enters Search mode to recover from degraded conditions. Fig.~\ref{fig:latency_breakdown} shows that the measured end-to-end runtime is composed of host control, transfer, accelerator execution, and synchronization overhead, with the Track path benefiting most strongly from feature reuse and mode-aware execution.

\begin{table}[!t]
\caption{Main Results of HBTXR}
\label{tab:main_results}
\centering
\footnotesize
\begin{tabular}{lccccc}
\toprule
Mode & Pupil Err. (px) & Gaze Err. (deg) & Search SR (\%) & Latency (ms) & Update Rate (Hz) \\
\midrule
Search & 0.58 & 0.91 & 97.8 & 0.93 & 1075 \\
Track & \textbf{0.42} & \textbf{0.68} & -- & \textbf{0.57} & \textbf{1754} \\
Scheduled Avg. & 0.45 & 0.71 & 97.8 & 0.63 & 1587 \\
\bottomrule
\end{tabular}
\end{table}

\begin{figure}[!t]
\centering
\draftfigure{fig6_qualitative.pdf}{0.95\columnwidth}{1.55in}
\caption{Qualitative tracking results of HBTXR under Search and Track modes. The examples show stable local updates in Track mode and robust recovery under degraded conditions through Search-mode re-entry.}
\label{fig:qualitative_results}
\end{figure}

\begin{figure}[!t]
\centering
\draftfigure{fig8_latency_breakdown.pdf}{0.95\columnwidth}{1.55in}
\caption{Latency breakdown of HBTXR. The total runtime is decomposed into host control, data transfer, accelerator execution, and synchronization overhead.}
\label{fig:latency_breakdown}
\end{figure}

\subsection{Comparison With Prior Eye-Tracking Methods}

Table~\ref{tab:compare_methods} summarizes representative prior works using the metrics reported in each source paper. Because these works differ in task definition, dataset, runtime platform, and evaluation protocol, the values are presented in their source-native form rather than forced into a single normalized ranking. This avoids mixing pupil-localization metrics, gaze-estimation metrics, and platform-specific runtime measurements into unsupported percentage comparisons.

The verified source values show that prior frame-based systems such as Etracker and HE-Tracker primarily report gaze accuracy and frame rate, whereas event-based systems such as E-Track and FACET report pupil-localization error and inference latency. EX-Gaze reports hybrid-system timing and EV-Eye pixel-error trade-offs under saccadic and smooth-pursuit settings. Accordingly, the most defensible cross-paper claim is that HBTXR combines explicit relocalization support with sub-millisecond tracking latency in a unified hybrid Search/Track framework.

The overall trend is visualized in Fig.~\ref{fig:tradeoff_plot}. The figure is intended as an illustrative qualitative summary rather than a strict one-to-one benchmark because the referenced works use heterogeneous metrics and deployment settings.

\begin{table*}[!t]
\caption{Verified Comparison With Prior Eye-Tracking Methods Using Source-Reported Metrics}
\label{tab:compare_methods}
\centering
\scriptsize
\begin{tabular}{l c l p{3.0cm} p{3.0cm} c p{2.2cm}}
\toprule
Work & Mod. & Primary Task & Reported Accuracy & Reported Runtime & Reloc. & Note \\
\midrule
Etracker & Frame & Gaze & 0.53$^\circ$ gaze accuracy & 30--60 Hz & Yes & Mobile near-eye gaze tracker \\
HE-Tracker & Frame & Gaze & 3.655$^\circ$ gaze error (all subjects) & 48 fps on AR HMD & Yes & Head-eye coordination model \\
E-Track & Event & Pupil & 3.68 px median center-distance error & 66.4 ms U-Net inference; 160 mW system & No & Event-only pupil tracking \\
FACET & Event & Pupil & 0.2030 px PE & 0.5302 ms inference & No & Event ellipse modeling \\
EX-Gaze & Hybrid & Pupil / tracking & 1.33 px (saccade), 1.49 px (smooth pursuit) on EV-Eye & 0.32 ms event tracking; 1.17 ms relocalization; $\sim$0.407 ms average op. & Yes & 2 KHz hybrid system \\
\textbf{HBTXR} & \textbf{Hybrid} & \textbf{Pupil + Gaze} & \textbf{0.42 px pupil err.; 0.68$^\circ$ gaze err.} & \textbf{0.57 ms Track mode; 1754 Hz} & \textbf{Yes} & \textbf{Unified Search/Track scheduler} \\
\bottomrule
\end{tabular}
\end{table*}

\begin{figure}[!t]
\centering
\draftfigure{fig7_tradeoff.pdf}{0.95\columnwidth}{1.55in}
\caption{Illustrative qualitative positioning of HBTXR and prior methods using source-reported metrics. Because the referenced works use heterogeneous tasks, datasets, and runtime definitions, the plot should be interpreted qualitatively rather than as a normalized ranking.}
\label{fig:tradeoff_plot}
\end{figure}
\subsection{Deployment-Oriented Analysis}

Table~\ref{tab:compare_deploy} reports deployment-oriented baselines in the metric form used by their source papers. We intentionally avoid converting energy-per-frame into power, or frame rate into latency, unless the source explicitly reports both. This prevents false equivalence across ASIC, FPGA SoC, and host-device systems.

The verified source values indicate that EyeCoD targets 240 FPS and reports 154.32 mW silicon power at 370 MHz, SEE-D reports 0.70 ms latency and 2.29 mJ per inference, and JaneEye reports 0.5 ms end-to-end latency, 2000 FPS, and 18.9 $\mu$J/frame. These baselines confirm that low-latency eye tracking is achievable in specialized deployments, but they also show that runtime metrics are highly platform- and metric-dependent. In this context, HBTXR is best positioned as a hybrid Search/Track system with runtime adaptivity rather than as a directly normalized replacement for any single event-only accelerator baseline.

This behavior is consistent with the mixed-granularity execution strategy shown in Fig.~\ref{fig:accel_overview} and Fig.~\ref{fig:search_track_timeline}. As emphasized in the captions of these figures, HBTXR uses phase-level execution for synchronization-heavy stages and tile-stream execution for reuse-friendly update stages, thereby reducing memory pressure and inter-stage idle gaps under realistic Search/Track operation.

\begin{table*}[!t]
\caption{Verified Deployment-Oriented Comparison Using Source-Reported Runtime and Efficiency Metrics}
\label{tab:compare_deploy}
\centering
\scriptsize
\begin{tabular}{l c c p{2.8cm} p{2.8cm} p{4.0cm}}
\toprule
System & Mod. & Runtime Adaptivity & Reported Runtime & Reported Efficiency & Deployment Scope \\
\midrule
EyeCoD & Frame & No & 240 FPS target throughput & 154.32 mW silicon power at 370 MHz; 335 mW config table & FlatCam-based frame pipeline \\
SEE-D & Event & Partial & 0.70 ms latency & 2.29 mJ/inference & Event-based FPGA SoC pipeline \\
JaneEye & Event & No & 0.5 ms; 2000 FPS & 18.9 $\mu$J/frame & Event-based 12-nm ASIC \\
\textbf{HBTXR} & \textbf{Hybrid} & \textbf{Yes} & \textbf{0.57 ms Track mode; 1754 Hz} & \textbf{Not reported} & \textbf{Hybrid Search/Track host-device deployment} \\
\bottomrule
\end{tabular}
\end{table*}
\subsection{Ablation Studies}

Table~\ref{tab:ablation} reports the ablation results. Removing the geometry-stable supervision increases the pupil-center error from 0.42 to 0.49 pixels and the gaze error from 0.68$^\circ$ to 0.80$^\circ$, confirming that the shared geometric target is critical for aligning Search and Track. Removing the decoupled head structure degrades both pupil and gaze accuracy to 0.47 pixels and 0.75$^\circ$, respectively, while also increasing latency to 0.62 ms.

The runtime-related ablations show a different trend. Disabling the Search/Track scheduler increases latency from 0.57 ms to 0.71 ms and degrades gaze accuracy to 0.73$^\circ$, showing that runtime adaptivity is important for both efficiency and stability. Disabling reuse-oriented execution leaves the accuracy unchanged at 0.42 pixels and 0.68$^\circ$, but increases latency from 0.57 ms to 0.69 ms, indicating that feature reuse is primarily a deployment-oriented gain rather than an accuracy mechanism.

Taken together, the ablations show that HBTXR derives its performance from three complementary sources: geometry-stable supervision for robust estimation, Search/Track decomposition for functional specialization, and reuse-oriented execution for low-latency deployment.

\begin{table}[!t]
\caption{Ablation Study of HBTXR}
\label{tab:ablation}
\centering
\footnotesize
\begin{tabular}{lccc}
\toprule
Variant & Pupil Err. (px) & Gaze Err. (deg) & Latency (ms) \\
\midrule
\textbf{Full HBTXR} & \textbf{0.42} & \textbf{0.68} & \textbf{0.57} \\
w/o geometry-stable supervision & 0.49 & 0.80 & 0.58 \\
w/o decoupled heads & 0.47 & 0.75 & 0.62 \\
w/o Search/Track scheduler & 0.46 & 0.73 & 0.71 \\
w/o reuse-oriented execution & 0.42 & 0.68 & 0.69 \\
\bottomrule
\end{tabular}
\end{table}

\subsection{Qualitative Analysis and Discussion}

The qualitative results in Fig.~\ref{fig:qualitative_results} show that HBTXR maintains smooth and stable tracking during normal operation and recovers effectively when confidence degrades. Meanwhile, the latency breakdown in Fig.~\ref{fig:latency_breakdown} indicates that the dominant practical benefit comes from reducing repeated refresh cost during steady-state tracking. This observation is consistent with the deployment hypothesis of HBTXR: the system benefits not only from a good predictor, but also from executing Search and Track differently.

Overall, the experiments support three conclusions. First, hybrid sensing alone is insufficient unless Search and Track are explicitly separated at the algorithmic and runtime levels. Second, geometry-stable supervision is critical for making frame-guided relocalization and event-guided residual tracking compatible within one model. Third, the deployment behavior of HBTXR is strongly improved by mode-aware mixed-granularity execution, which lowers the cost of steady-state tracking while preserving robust relocalization capability.

\section{Conclusion}

This paper presented \textbf{HBTXR}, an algorithm-hardware co-designed hybrid eye-tracking system for on-device XR applications. The proposed system was motivated by the observation that practical XR eye tracking cannot be addressed adequately by frame-only or event-only pipelines alone. Frame-based methods provide semantically stable appearance cues but incur high bandwidth and latency costs, whereas event-based methods offer fast and sparse sensing but remain weaker in relocalization and long-horizon stability. To bridge this gap, HBTXR formulates eye tracking as a hybrid Search-and-Track problem and jointly considers supervision, model design, runtime scheduling, and deployment architecture within a unified system framework.

The proposed approach consists of four tightly connected components. First, HBTXR introduces a geometry-stable hybrid supervision framework that stabilizes geometric targets across frame and event inputs. Second, it employs a shared-backbone hybrid transformer with decoupled Search, Event, Track, and Mask heads to support both robust relocalization and low-latency residual update. Third, it incorporates a runtime-aware Search/Track scheduler that explicitly switches between refresh-oriented and state-persistent operating modes according to confidence, similarity, event density, and eye-state signals. Fourth, it presents a mode-aware mixed-granularity deployment architecture, enabling the hybrid tracking algorithm to be mapped onto a shared accelerator substrate with feature reuse and low-overhead control.

Experimental results demonstrated that HBTXR achieves a pupil-center error of \textbf{0.42 pixels}, a gaze error of \textbf{0.68$^\circ$}, and an end-to-end latency of \textbf{0.57 ms} in Track mode. These results confirm that HBTXR provides a practical balance among semantic robustness, temporal responsiveness, and deployment-oriented efficiency. More importantly, the results show that the main value of HBTXR lies in its integrated operating point: frame-guided Search improves recovery robustness, event-guided Track improves low-latency update, and the scheduler together with the deployment-aware execution path makes the combined system practical for on-device XR scenarios.

Several directions remain open. The Search path still incurs higher cost than steady-state Track operation and can be further optimized for more aggressive runtime efficiency. The hardware realization can also be extended beyond the current deployment-oriented architecture toward deeper implementation and energy/resource measurement. Finally, the proposed framework may be generalized to broader eye-related tasks such as blink understanding, user-state estimation, and other multimodal XR interaction primitives.

In summary, HBTXR demonstrates that effective on-device XR eye tracking requires more than a high-accuracy model or a fast sensing path in isolation. Instead, robust supervision, hybrid Search/Track formulation, runtime-aware scheduling, and deployment-aware execution must be designed together. The proposed system provides a concrete step toward such an integrated eye-tracking architecture for next-generation XR devices.





\bibliographystyle{IEEEtran}
\bibliography{ieee.bib, main.bib}

\end{document}



